\documentclass[runningheads]{llncs}

\usepackage{graphicx}
\usepackage{amsmath}

\begin{document}

\title{Experimental Analysis of Graph Coloring Algorithms}

\author{
Albert-Iulian Romaniuc \and
C\u{a}lin-Andrei Necula \and
\c{S}tefan Solomon
}

\institute{Proiect - Analiza Algoritmilor}

\maketitle

\begin{abstract}
Colorarea grafurilor este o problemă clasică de optimizare combinatorială,
cunoscută ca fiind NP-hard, cu aplicații în planificare, alocarea
registrelor și analiza rețelelor. În această lucrare studiem problema
NP-2 Graph Coloring prin compararea unei soluții exacte bazate pe backtracking
cu doi algoritmi euristici bine cunoscuți: DSATUR și Recursive Largest First
(RLF). Algoritmul exact este utilizat ca referință pentru evaluarea calității
soluțiilor euristice. Analiza experimentală este realizată pe grafuri generate
aleator, de dimensiuni și densități diferite, luând în considerare atât timpul
de execuție, cât și numărul de culori utilizate. Rezultatele evidențiază
compromisul dintre optimalitate și eficiență și arată că metodele euristice
pot furniza soluții apropiate de optim într-un timp de calcul mult mai redus.
\end{abstract}

\section{Introducere}

Problema colorării grafurilor este una dintre problemele fundamentale din
teoria grafurilor și optimizarea combinatorială. Dat fiind un graf neorientat,
scopul este atribuirea unei culori fiecărui vârf astfel încât două vârfuri
adiacente să nu aibă aceeași culoare, utilizând un număr minim de culori.
Determinarea numărului cromatic al unui graf este o problemă NP-hard,
ceea ce face ca soluțiile exacte să devină rapid imposibil de calculat pentru
instanțe de dimensiuni mari.

Din cauza acestei complexități, au fost propuse numeroase algoritmi euristici
și aproximativi care urmăresc obținerea unor soluții de bună calitate într-un
timp rezonabil. Printre aceștia, algoritmii DSATUR și Recursive Largest First
(RLF) sunt metode greedy populare, utilizate frecvent în practică datorită
echilibrului dintre calitatea soluției și costul computațional. Cu toate
acestea, performanța algoritmilor euristici depinde puternic de structura
instanței analizate și este dificil de evaluat fără o soluție de referință.

În această lucrare realizăm o analiză experimentală comparativă între un
algoritm exact de tip backtracking și două abordări euristice, DSATUR și RLF.
Algoritmul exact este folosit ca reper pentru evaluarea calității soluțiilor
euristice, folosind numărul de culori ca metrică principală. Performanța
algoritmilor este analizată atât din punct de vedere al timpului de execuție,
cât și al calității soluției, pe grafuri cu dimensiuni și densități variate.
Scopul studiului este evidențierea compromisurilor dintre optimalitate și
eficiență în rezolvarea problemei NP-2 Graph Coloring.

\section{Definirea problemei}

Fie $G = (V, E)$ un graf neorientat, unde $V$ reprezintă mulțimea vârfurilor,
iar $E \subseteq V \times V$ reprezintă mulțimea muchiilor. O colorare a
grafului este o funcție $c : V \rightarrow \mathbb{N}$ care asociază fiecărui
vârf o culoare, astfel încât pentru orice muchie $(u, v) \in E$ să fie
îndeplinită condiția $c(u) \neq c(v)$.

Numărul total de culori distincte utilizate într-o colorare este cardinalul
mulțimii $\{c(v) \mid v \in V\}$. O colorare este considerată validă dacă
respectă restricția de adiacență, iar obiectivul problemei este minimizarea
numărului de culori folosite.

Numărul cromatic al grafului $G$, notat cu $\chi(G)$, este definit ca fiind
numărul minim de culori necesare pentru a obține o colorare validă a grafului.
Determinarea exactă a numărului cromatic este o problemă NP-hard, ceea ce
implică faptul că nu există algoritmi cunoscuți care să rezolve problema în
timp polinomial pentru toate instanțele.

În cadrul acestei lucrări, ne concentrăm asupra problemei NP-2 Graph Coloring,
care constă în determinarea unei colorări valide cu un număr minim de culori,
folosind atât metode exacte, cât și euristice, și în compararea performanței
acestora din punct de vedere al calității soluției și al timpului de execuție.

\section{Algoritmi}

În această secțiune sunt prezentate metodele utilizate pentru rezolvarea
problemei de colorare a grafurilor. Sunt analizate atât o abordare exactă,
bazată pe backtracking, cât și două metode euristice, DSATUR și Recursive
Largest First (RLF), care oferă soluții aproximative într-un timp de execuție
semnificativ mai redus.

\subsection{Algoritmul de backtracking}

Algoritmul de backtracking explorează în mod sistematic toate posibilitățile
de atribuire a culorilor vârfurilor unui graf, construind soluția pas cu pas.
La fiecare etapă, se verifică dacă atribuirea curentă respectă restricțiile
problemei, iar ramurile care conduc la soluții invalide sunt eliminate prin
procesul de pruning.

Această abordare garantează obținerea unei soluții optime, deoarece toate
colorările posibile sunt analizate. Totuși, complexitatea sa este exponențială
în raport cu numărul de vârfuri, ceea ce limitează aplicabilitatea algoritmului
la grafuri de dimensiuni reduse.

\subsection{Algoritmul DSATUR}

DSATUR (Degree of Saturation) este un algoritm euristic greedy care selectează
pentru colorare, la fiecare pas, vârful cu gradul de saturație maxim. Gradul
de saturație al unui vârf reprezintă numărul de culori distincte utilizate de
vecinii săi deja colorați.

Prin prioritizarea vârfurilor cu cele mai restrictive condiții, DSATUR tinde
să reducă numărul total de culori utilizate. Algoritmul oferă în practică
soluții de calitate apropiată de optim, având un timp de execuție redus, însă
nu garantează obținerea soluției optime pentru toate instanțele.

\subsection{Algoritmul Recursive Largest First (RLF)}

Algoritmul Recursive Largest First construiește colorarea grafului prin
identificarea succesivă a unor mulțimi independente de vârfuri, cărora li se
atribuie aceeași culoare. La fiecare iterație, algoritmul selectează un vârf
cu grad mare și extinde mulțimea independentă prin adăugarea vârfurilor care
nu sunt adiacente cu cele deja selectate.

RLF este deosebit de eficient pentru grafuri dense, unde identificarea unor
mulțimi independente mari poate conduce la o reducere a numărului total de
culori. Cu toate acestea, performanța sa depinde de structura grafului, iar
calitatea soluției poate varia de la o instanță la alta.

\section{Setarea experimentală}

Analiza experimentală are ca scop compararea performanței algoritmilor
prezentați din punct de vedere al timpului de execuție și al calității
soluției obținute. Pentru a permite utilizarea algoritmului exact de tip
backtracking ca soluție de referință, experimentele au fost limitate la
grafuri cu un număr redus de vârfuri.

\subsection{Generarea instanțelor de test}

Instanțele de test sunt grafuri neorientate generate aleator, cu un număr de
vârfuri cuprins între 4 și 20. Această limitare este necesară deoarece
complexitatea exponențială a algoritmului de backtracking face imposibilă
rezolvarea instanțelor de dimensiuni mai mari într-un timp rezonabil.

Pentru fiecare dimensiune a grafului au fost generate instanțe cu densități
diferite, definite ca raportul dintre numărul de muchii existente și numărul
maxim posibil de muchii. Prin variarea densității se urmărește analiza
influenței structurii grafului asupra performanței algoritmilor.

\subsection{Metrici de evaluare}

Performanța algoritmilor este evaluată folosind două metrici principale:
timpul de execuție și numărul de culori utilizate în soluția finală. Timpul
de execuție este măsurat în secunde și reflectă eficiența algoritmului, în
timp ce numărul de culori utilizate reprezintă calitatea soluției obținute.

Algoritmul de backtracking este utilizat pentru determinarea soluției optime,
iar rezultatele algoritmilor euristici sunt comparate direct cu această
referință. În acest mod se poate evalua cât de apropiate sunt soluțiile
euristice față de optim.

\subsection{Mediul de rulare}

Toate experimentele au fost realizate în aceleași condiții de rulare, folosind
aceleași instanțe de test pentru fiecare algoritm, pentru a asigura o
comparație corectă. Algoritmii au fost implementați în limbajul C, iar
rularea și colectarea rezultatelor au fost automatizate prin scripturi
auxiliare.

\section{Rezultate experimentale}

\subsection{Timpul de execuție}

\subsubsection{Backtracking}

Algoritmul de backtracking prezintă o creștere rapidă a timpului de execuție
odată cu mărirea dimensiunii grafului. Acest comportament este așteptat, deoarece
algoritmul explorează exhaustiv spațiul soluțiilor posibile.

\begin{figure}[h]
\centering
\includegraphics[width=0.85\linewidth]{backtracking_time.png}
\caption{Timpul de execuție al algoritmului de backtracking.}
\label{fig:bt_time}
\end{figure}

Chiar și pentru grafuri cu un număr relativ mic de vârfuri, variațiile timpului
de execuție devin vizibile, confirmând complexitatea exponențială a metodei.
Acest fapt limitează utilizarea practică a algoritmului la instanțe de
dimensiuni reduse.

\newpage

\subsubsection{DSATUR}

Algoritmul prezintă un timp de execuție stabil pentru toate dimensiunile
analizate. Spre deosebire de algoritmul de backtracking, care manifestă o
creștere rapidă și fluctuații semnificative ale timpului de rulare, acesta
menține valori reduse și relativ constante ale timpului de execuție pe întreg
intervalul de grafuri testate. Acest comportament indică o scalabilitate
superioară și o sensibilitate redusă la creșterea dimensiunii instanței.

\begin{figure}[h]
\centering
\includegraphics[width=0.85\linewidth]{dsatur_time.png}
\caption{Timpul de execuție al algoritmului DSATUR.}
\label{fig:dsatur_time}
\end{figure}

Stabilitatea observată în Figura~\ref{fig:dsatur_time} este o consecință directă
a caracterului euristic al algoritmului. DSATUR evită explorarea exhaustivă a
spațiului soluțiilor și construiește colorarea incremental, selectând la fiecare
pas vârful cu cel mai mare grad de saturație, adică vârful adiacent celor mai
multe culori distincte deja utilizate.

Prin prioritizarea vârfurilor cu cele mai restrictive condiții, algoritmul
reduce probabilitatea apariției conflictelor în etapele târzii ale procesului
de colorare. Această strategie limitează numărul de decizii costisitoare și
conduce la un timp de execuție predictibil, chiar și pentru grafuri cu un număr
mai mare de vârfuri.

\newpage

\subsubsection{Recursive Largest First}

Algoritmul Recursive Largest First (RLF) prezintă un timp de execuție redus,
comparabil cu cel al algoritmului DSATUR, însă cu variații mai pronunțate pe
anumite instanțe. Spre deosebire de DSATUR, care manifestă un comportament
temporal foarte stabil, RLF este mai sensibil la structura grafului analizat,
în special la densitatea acestuia și la distribuția muchiilor.

\begin{figure}[h]
\centering
\includegraphics[width=0.85\linewidth]{rlf_time.png}
\caption{Timpul de execuție al algoritmului RLF.}
\label{fig:rlf_time}
\end{figure}

Fluctuațiile observate în Figura~\ref{fig:rlf_time} sunt cauzate de modul de
funcționare al algoritmului, care construiește colorarea prin identificarea
succesivă a unor mulțimi independente maximale. Procesul de extindere a acestor
mulțimi poate necesita un efort computațional suplimentar, în special pentru
grafuri dense, unde restricțiile de adiacență sunt mai numeroase.

În consecință, timpul de execuție al algoritmului RLF poate varia în funcție de
structura instanței, chiar dacă dimensiunea grafului rămâne constantă. Cu toate
acestea, algoritmul rămâne eficient din punct de vedere computațional și oferă
performanțe competitive pentru o gamă largă de instanțe.

\newpage

\subsubsection{Comparație între algoritmi}

Figura~\ref{fig:time_vs_vertices} prezintă o comparație directă a timpului de
execuție pentru cei trei algoritmi analizați — backtracking, DSATUR și Recursive
Largest First — în funcție de dimensiunea grafului. Se observă diferențe clare
între abordarea exactă și metodele euristice. Algoritmul de backtracking
manifestă o creștere rapidă a timpului de execuție odată cu mărirea numărului de
vârfuri, comportament determinat de explorarea exhaustivă a spațiului
soluțiilor, caracteristică problemelor NP-hard. Chiar și pentru grafuri de
dimensiuni relativ reduse, timpul necesar devine semnificativ, limitând
aplicabilitatea practică a metodei.

În contrast, algoritmii euristici DSATUR și RLF prezintă timpi de execuție mult
mai reduși și o evoluție controlată pe măsură ce dimensiunea instanței crește.
DSATUR se remarcă printr-un comportament deosebit de stabil, cu variații minime,
datorită strategiei sale bazate pe gradul de saturație, care prioritizează
vârfurile critice încă din etapele incipiente ale colorării. Algoritmul RLF
prezintă timpi comparabili, însă cu fluctuații mai vizibile, influențate de
structura grafului și de dificultatea identificării mulțimilor independente.
Per ansamblu, comparația evidențiază compromisul fundamental dintre
optimalitate și eficiență: algoritmul exact garantează soluția optimă, dar este
impracticabil pentru instanțe mai mari, în timp ce metodele euristice oferă
soluții eficiente într-un timp de calcul semnificativ mai redus.

\begin{figure}[h]
\centering
\includegraphics[width=0.9\linewidth]{time_vs_vertices.png}
\caption{Comparația timpului de execuție între algoritmi.}
\label{fig:time_vs_vertices}
\end{figure}

\newpage

\subsection{Numărul de culori utilizate}

\subsubsection{Backtracking}

Algoritmul de backtracking furnizează în toate cazurile soluția optimă,
reprezentând un reper pentru evaluarea performanței algoritmilor euristici.
Numărul de culori utilizate reflectă direct numărul cromatic al grafului
analizat și este independent de ordinea de procesare a vârfurilor, deoarece
algoritmul explorează exhaustiv toate colorările valide posibile.

\begin{figure}[h]
\centering
\includegraphics[width=0.85\linewidth]{backtracking_colors.png}
\caption{Numărul de culori utilizate de algoritmul de backtracking în funcție de densitatea grafului.}
\label{fig:bt_colors}
\end{figure}

Se observă că, pe măsură ce densitatea grafului crește, numărul de culori
necesare crește de asemenea. Acest comportament este așteptat, deoarece grafurile
mai dense impun mai multe restricții de adiacență între vârfuri, reducând
posibilitatea reutilizării aceleiași culori pentru vârfuri neadiacente.

În consecință, creșterea densității conduce la o creștere a numărului cromatic
al grafului, iar algoritmul de backtracking reflectă fidel această evoluție,
furnizând valoarea optimă pentru fiecare instanță analizată. Din acest motiv,
rezultatele obținute prin backtracking sunt utilizate ca referință în analiza
comparativă a calității soluțiilor euristice.

\newpage

\subsubsection{DSATUR}

Algoritmul DSATUR urmărește obținerea unei colorări eficiente prin
prioritizarea vârfurilor cu grad mare de saturație, ceea ce conduce, în
practică, la soluții apropiate de optim din punct de vedere al numărului de
culori utilizate. Prin selectarea vârfurilor care au cele mai multe restricții
impuse de vecinătatea deja colorată, algoritmul limitează apariția conflictelor
și reduce necesitatea introducerii unor culori suplimentare.

\begin{figure}[h]
\centering
\includegraphics[width=0.85\linewidth]{dsatur_colors.png}
\caption{Numărul de culori utilizate de algoritmul DSATUR în funcție de densitatea grafului.}
\label{fig:dsatur_colors}
\end{figure}

Rezultatele arată că DSATUR obține, în majoritatea instanțelor testate, un număr
de culori foarte apropiat de soluția optimă furnizată de backtracking.
Comportamentul algoritmului este relativ stabil pentru densități diferite, iar
abaterile față de optim sunt reduse, chiar și pe măsură ce graful devine mai
dens.

Această stabilitate indică faptul că strategia bazată pe gradul de saturație
reprezintă un criteriu eficient de selecție a vârfurilor în procesul de
colorare. În consecință, DSATUR constituie o alternativă practică la metodele
exacte pentru grafuri de dimensiuni mai mari, oferind un compromis favorabil
între calitatea soluției și costul computațional.

\newpage

\subsubsection{Recursive Largest First}

Algoritmul Recursive Largest First construiește colorarea prin identificarea
succesivă a unor mulțimi independente, iar calitatea soluției obținute depinde
în mod direct de modul în care aceste mulțimi pot fi extinse în funcție de
structura grafului. Strategia urmărește atribuirea aceleiași culori unui număr
cât mai mare de vârfuri neadiacente, însă eficiența acestei abordări este
influențată de densitatea și distribuția muchiilor.

\begin{figure}[h]
\centering
\includegraphics[width=0.85\linewidth]{rlf_colors.png}
\caption{Numărul de culori utilizate de algoritmul RLF în funcție de densitatea grafului.}
\label{fig:rlf_colors}
\end{figure}

Rezultatele indică o variabilitate mai mare a numărului de culori utilizate
comparativ cu DSATUR. Pentru grafuri dense, identificarea unor mulțimi
independente mari devine mai dificilă, ceea ce limitează posibilitatea de a
colora simultan mai multe vârfuri cu aceeași culoare și conduce, în consecință,
la utilizarea unui număr mai mare de culori.

Cu toate acestea, în anumite instanțe, algoritmul RLF oferă rezultate
competitive, apropiate de optim, în special atunci când structura grafului
permite formarea unor mulțimi independente suficient de mari. Acest comportament
confirmă faptul că performanța algoritmului este puternic dependentă de
structura instanței analizate.

\newpage

\subsubsection{Comparație între algoritmi}

Figura~\ref{fig:colors_vs_density} prezintă o comparație directă a numărului de
culori utilizate de cei trei algoritmi analizați — backtracking, DSATUR și
Recursive Largest First — în funcție de densitatea grafului. Algoritmul de
backtracking este utilizat ca referință, deoarece furnizează în toate cazurile
soluția optimă, stabilind valoarea exactă a numărului cromatic. Algoritmii
euristici DSATUR și RLF tind să aproximeze această valoare, însă cu grade diferite
de stabilitate și acuratețe.

Experimentele au fost rulate pe aceeași platformă hardware, utilizând un procesor
Intel Core i5, cu arhitectură \emph{x86\_64}, și 8\,GB de memorie RAM. Execuția a
fost realizată în regim single-thread, pentru a elimina influența paralelismului
asupra rezultatelor. Toți algoritmii au fost compilați și executați în aceleași
condiții de rulare, pe același sistem de operare și cu aceleași opțiuni de
optimizare ale compilatorului, astfel încât diferențele observate să reflecte
exclusiv comportamentul algoritmic și strategiile de colorare utilizate.

Rezultatele arată că DSATUR se remarcă printr-un comportament constant și soluții
foarte apropiate de optim pentru majoritatea instanțelor, indiferent de
densitatea grafului. Algoritmul RLF prezintă fluctuații mai pronunțate, ceea ce
indică o sensibilitate mai mare la structura grafului și la distribuția
muchilor. În ansamblu, comparația evidențiază faptul că, deși metodele euristice
nu garantează obținerea soluției optime, acestea pot furniza soluții de bună
calitate într-un cadru experimental controlat, cu un compromis acceptabil față
de soluția optimă.

\begin{figure}[h]
\centering
\includegraphics[width=0.9\linewidth]{colors_vs_density.png}
\caption{Comparația numărului de culori utilizate în funcție de densitatea grafului.}
\label{fig:colors_vs_density}
\end{figure}

\newpage

\section{Concluzii}

În această lucrare a fost realizată o analiză experimentală a problemei
NP-2 Graph Coloring, comparând un algoritm exact bazat pe backtracking cu două
metode euristice, DSATUR și Recursive Largest First. Rezultatele obținute
confirmă faptul că algoritmul de backtracking garantează soluția optimă, însă
aplicabilitatea sa este limitată la grafuri de dimensiuni reduse din cauza
complexității exponențiale.

Algoritmii euristici analizați oferă soluții apropiate de optim într-un timp
de execuție semnificativ mai mic, fiind mai potriviți pentru instanțe de
dimensiuni mai mari. DSATUR se remarcă printr-un comportament stabil și un
echilibru bun între calitatea soluției și eficiență, în timp ce RLF poate fi
avantajos pentru grafuri dense, dar cu o variabilitate mai mare a rezultatelor.

În concluzie, alegerea algoritmului potrivit depinde de dimensiunea și
structura grafului analizat. Metodele exacte sunt utile ca referință teoretică,
în timp ce euristicile reprezintă soluții practice eficiente pentru aplicații
reale. Studiul evidențiază compromisurile fundamentale dintre optimalitate și
eficiență în rezolvarea problemelor NP-hard de optimizare combinatorială.

\begin{thebibliography}{8}

\bibitem{garey}
M. R. Garey and D. S. Johnson,
\textit{Computers and Intractability: A Guide to the Theory of NP-Completeness},
W. H. Freeman, 1979.

\bibitem{dsatur}
D. Brélaz,
New methods to color the vertices of a graph,
\textit{Communications of the ACM}, vol. 22, no. 4, pp. 251--256, 1979.

\bibitem{rlf}
F. T. Leighton,
A graph coloring algorithm for large scheduling problems,
\textit{Journal of Research of the National Bureau of Standards},
vol. 84, pp. 489--506, 1979.

\bibitem{graphcoloring}
I. Jensen and B. Toft,
\textit{Graph Coloring Problems},
Wiley-Interscience, 1995.

\bibitem{heuristics}
T. F. Gonzalez,
\textit{Handbook of Approximation Algorithms and Metaheuristics},
Chapman \& Hall/CRC, 2007.

\end{thebibliography}

\end{document}